% generated by build_assets.py -- do not edit
We describe and evaluate LEBRE, an online one-step-ahead forecaster for a target driven by exogenous inputs, at a cost of a few hundred floating-point operations per step. The live model is a memory of the target's own past and an online linear regression on the current inputs, combined by dynamic model averaging. Structural changes (adding, removing or swapping a whole input, or refining its lag response) run in shadow and are accepted only when an anytime-valid e-process of predictive improvement crosses a level set by an online multiple-testing rule. The components are not new; the contribution is their integration at a nearly fixed cost, a condition under which predictive improvement may be read as structure, and a pre-registered evaluation. When the reference model is learned online, its estimation error is predictable, so a challenger can improve on it without new structure; the structural reading requires small reference misadjustment. In three sequential pre-registered evaluations on never-seen hydrology and building-energy series, the first (125 series) revealed an engineering failure. After the fix, on 40 and 60 series, LEBRE had the lowest error among the low-cost comparators evaluated (MSE relative to online NLinear 0.812 on 60 series, 95\% CI 0.752--0.866) and no catastrophic series; it tied a per-series SARIMAX with inputs and had about 25\% higher error than Chronos-2 with covariates. The pre-declared cost contract was narrowly missed (+3.7\% mean, +2.0\% peak). A C99 port ran on a simulated Cortex-M4F at 3.2--6.0 thousand instructions per step with 73 KB of state. The testing mechanism did not improve accuracy over switching all inputs on. Code, pre-registrations and per-series result tables are public.
